\chapter{Lebesgue measure}
\lecture{8}{19 March.}{}
To overcome the limitations of the Riemann integral, particularly its inability to handle pointwise limits of functions, this chapter constructs a more robust theory of integration starting with the concept of measure.

\section{The goal: Lebsgue measure}

Our goal is to define a concept of measurable set, which will be a special kind of subset of \(\mathbb{R} ^n\). For every such measurable set \(\Omega \subseteq \mathbb{R} ^n\), we'll define the Lebesgue measure \(m(\Omega )\) to be a number in \([0, \infty]\).
\begin{definition}[Measurable set]
    The concept of measurable set will obey the following properties:
    \begin{itemize}
        \item [(i)] \textbf{(Borel property)} Every open set in \(\mathbb{R} ^n\) is measurable, as is every closed set. 
        \item [(ii)] \textbf{(Complementarity)} If \(\Omega \) is measurable, then \(\mathbb{R} ^n \setminus \Omega \) is also measurable.
        \item [(iii)] \textbf{(Boolean algebra propery)} If \((\Omega_j )_{j \in J}\) is any finite collection of measurable sets, then \(\bigcup_{j \in J} \Omega _j \) and \(\bigcap_{j \in J} \Omega _j \) are also measurable. 
        \item [(iv)] \textbf{(\(\sigma \)-algebra property)} If \((\Omega _j)_{j \in J}\) are any countable collection of measurable sets (so \(J\) is countable), then \(\bigcup_{j \in J} \Omega _j \) and \(\bigcap_{j \in J} \Omega _j \) are also measurable.       
    \end{itemize}
\end{definition}

Note that some of these properties are redundant; for instance, (iv) will imply (iii), and once one knows all open sets are measurable, (ii) will imply that all closed sets are measurable also. The properties (i)-(iv) will ensure that virtually every set one cares about is measurable; though there do exist non-measurable sets.

\begin{definition}[Lebsegue measure]
    To every measurable set \(\Omega \), we associate the Lebsegue measure \(m(\Omega )\) of \(\Omega \), which will obey the following properties:
    \begin{itemize}
        \item [(v)] \textbf{(Empty set)} The empty set \(\varnothing \) has measure \(m(\varnothing ) = 0\).  
        \item [(vi)] \textbf{(Positivity)} We have \(0 \le m(\Omega ) \le +\infty \) for every measurable set \(\Omega \).   
        \item [(vii)] \textbf{(Monotocity)} If \(A \subseteq B\), and \(A\) and \(B\) are both measurable, then \(m(A) \le m(B)\).     
        \item [(viii)] \textbf{(Finite sub-additivity)} If \((A_j)_{j \in J}\) are a finite collection of measurable sets, then 
        \[
            m \left( \bigcup_{j \in J} A_j  \right) \le \sum_{j \in J} m(A_j).  
        \]
        \item [(ix)] \textbf{(Finite additivity)} If \((A_j)_{j \in J}\) are a finite collection of \textit{disjoint} measurable sets, then 
        \[
            m \left( \bigcup_{j \in J} A_j  \right) = \sum_{j \in J} m(A_j).  
        \]
        \item [(x)] \textbf{(Countable sub-additivity)} If \((A_j)_{j \in J}\) are a countable collection of measurable sets, then 
        \[
            m\left( \bigcup_{j \in J} A_j  \right) \le \sum_{j \in J} m(A_j).  
        \]
        \item [(xi)] \textbf{(Countable additivity)} If \((A_j)_{j \in J}\) are a countable collection of \textit{disjoint} measurable sets, then 
        \[
            m \left( \bigcup_{j \in J} A_j  \right) = \sum_{j \in J} m(A_j).  
        \] 
        \item [(xii)] \textbf{(Normalization)} The unit cube
        \[
            [0, 1]^n = \left\{ (x_1, \dots , x_n) \in \mathbb{R} ^n: 0 \le x_j \le 1 \text{ for all } 1 \le j \le n  \right\} 
        \]
        has measure \(m([0, 1]^n) = 1\). 
        \item [(xiii)] \textbf{(Translation invariance)} If \(\Omega \) is a measurable set, and \(x \in \mathbb{R} ^n\), then \(x + \Omega \coloneqq \left\{ x + y: y \in \Omega  \right\} \) is also measurable, and \(m(x + \Omega ) = m (\Omega )\).     
    \end{itemize}  
\end{definition}

Again, many of these properties are redundant; for instance, the countable additivity property can be used to deduce the finite additivity property, which in turn can be used to derive monotonicity (when combined with the positivity property). One can also obtain the sub-additivity properties from the additivity ones. Note that \(m(\Omega )\) can be \(+ \infty \), and so in particular some of the sums in the above properties may also equal \(+\infty \); in this chapter we adopt the convention that an infinite sum \(\sum_{j \in J} a_j \) of non-negative quantities \(a_j\) is equal to \(+ \infty \) if the sum is not absolutely convergent. (Since everything is non0negative we will never have to deal with indeterminate forms such as \((- \infty) + (+ \infty ) \).)

Our goal for this chapter can then be stated:

\begin{theorem}[Existence of Lebsegue measure]
    There exists a concept of a measurable set and a way to assign a number \(m(\Omega )\) to every measurable set \(\Omega \subseteq \mathbb{R} ^n\), which obeys all of the properties (i) to (xiii).  
\end{theorem}